\chapter{Tests et tests multiples}\label{ch:qm:testing-and-multiple-testing}

Une équipe de recherche teste deux cents variantes d'un signal de momentum et présente la meilleure~: sa $t$-statistique de backtest vaut 3.1,
une $p$-valeur unilatérale de un pour mille. Si aucune des variantes n'avait d'edge, et si elles avaient été indépendantes, la plus grande de
deux cents $t$-statistiques dépasserait encore 3.1 avec une probabilité de 17.6~\%, soit environ une fois sur six. Les variantes ne sont pas
indépendantes, ce qui aide~; l'équipe a aussi choisi la période d'échantillon, l'univers et le modèle de coûts, ce qui, lui, n'aide pas. Ce
chapitre porte sur ce qu'un test peut établir~: le niveau et la puissance, et le lemme qui fait du rapport de vraisemblance la bonne statistique~; les
nombreux choix discrets qui multiplient les tests qu'un chercheur effectue~; les corrections qui contrôlent la probabilité d'une fausse découverte,
ou leur proportion~; et les tests conçus pour la recherche de stratégies, qui tiennent compte de la corrélation entre elles.

\section{Tests, niveau et puissance}

\begin{definition}[Test d'hypothèse, hypothèse nulle, $p$-valeur]\label{def:qm:testing-and-multiple-testing:test}
Un \emph{test d'hypothèse}\index{test d'hypothèse} est une règle qui, à partir des données $X$, décide s'il faut rejeter une \emph{hypothèse
nulle}\index{hypothèse nulle} $H_0$ (une affirmation sur la loi de $X$, telle que «~l'espérance de rendement de la stratégie est nulle~») au profit
d'une alternative $H_1$. Il rejette lorsqu'une statistique $T(X)$ dépasse une valeur critique. La \emph{$p$-valeur}\index{p-valeur} d'un
$t$ observé est $\P_{H_0}(T \ge t)$~: la probabilité, sous l'hypothèse nulle, d'une statistique au moins aussi extrême que celle observée.
\end{definition}

Une $p$-valeur n'est pas la probabilité que l'\omterm{def:qm:testing-and-multiple-testing:test}{hypothèse nulle} soit vraie, et un moins cette valeur n'est pas la probabilité que la stratégie fonctionne~:
ces deux quantités exigent une loi a priori (chapitre~14). Elle constitue une affirmation sur les données sous une seule hypothèse.

\begin{definition}[Niveau, puissance]\label{def:qm:testing-and-multiple-testing:size}
Le \emph{niveau d'un test}\index{niveau d'un test} est la plus grande probabilité de rejeter lorsque l'\omterm{def:qm:testing-and-multiple-testing:test}{hypothèse nulle} est vraie (un faux
positif, ou erreur de type~I). La \emph{puissance d'un test}\index{puissance d'un test} en une alternative donnée est la probabilité de rejeter
lorsque cette alternative est vraie~; un moins la puissance est l'erreur de type~II.
\end{definition}

\begin{theorem}[Lemme de Neyman--Pearson]\label{thm:qm:testing-and-multiple-testing:np}
Pour une \omterm{def:qm:testing-and-multiple-testing:test}{hypothèse nulle} simple (densité $f_0$) contre une alternative simple ($f_1$), le test qui rejette lorsque $f_1(x)/f_0(x) > k$, avec $k$
choisi pour que son niveau soit $\alpha$, est uniformément le plus puissant~: aucun test de niveau au plus $\alpha$ n'a une puissance supérieure.
\end{theorem}

\begin{proof}
Soit $\phi^*$ le \omterm{def:qm:testing-and-multiple-testing:lr}{test du rapport de vraisemblance} et $\phi$ un test quelconque de niveau au plus $\alpha$ (toutes deux fonctions indicatrices de la
région de rejet). Point par point $(\phi^* - \phi)(f_1 - kf_0) \ge 0$~: là où $\phi^* = 1$, $f_1 > kf_0$~; là où $\phi^* = 0$, $f_1 \le kf_0$.
En intégrant, $\text{puissance}(\phi^*) - \text{puissance}(\phi) \ge k(\alpha - \text{niveau}(\phi)) \ge 0$.
\end{proof}

Pour la moyenne de rendements normaux, le rapport $f_1/f_0$ contre une espérance positive quelconque est croissant en la moyenne empirique, de sorte que le
$t$-test unilatéral est simultanément le test le plus puissant de «~pas d'edge~» contre tout edge positif. En présence d'hypothèses composites et
de paramètres de nuisance, le lemme devient une recette.

\begin{definition}[Test du rapport de vraisemblance]\label{def:qm:testing-and-multiple-testing:lr}
Le \emph{test du rapport de vraisemblance}\index{test du rapport de vraisemblance} d'une \omterm{def:qm:testing-and-multiple-testing:test}{hypothèse nulle} qui impose $r$ restrictions à $\theta$ rejette pour de grandes valeurs de
$\mathrm{LR} = 2\bigl(\ell(\hat\theta) - \ell(\hat\theta_0)\bigr)$, où $\ell$ est la log-vraisemblance, $\hat\theta$ son maximiseur et
$\hat\theta_0$ son maximiseur sous les restrictions.
\end{definition}

D'après le théorème de Wilks, sous l'\omterm{def:qm:testing-and-multiple-testing:test}{hypothèse nulle} et des conditions de régularité (le vrai paramètre appartenant à l'espace des paramètres), $\mathrm{LR}$ est
asymptotiquement une loi $\chi^2_r$. La condition compte en pratique~: tester qu'une excitation de Hawkes (chapitre~7) est nulle place
l'\omterm{def:qm:testing-and-multiple-testing:test}{hypothèse nulle} au bord $\alpha = 0$, et la limite devient un mélange égal d'une masse ponctuelle en zéro et d'une loi $\chi^2_1$, ce qui divise
la $p$-valeur par deux.

\begin{definition}[Test de Kolmogorov--Smirnov]\label{def:qm:testing-and-multiple-testing:ks}
Le \emph{test de Kolmogorov--Smirnov}\index{test de Kolmogorov--Smirnov} de l'\omterm{def:qm:testing-and-multiple-testing:test}{hypothèse nulle} selon laquelle $X_1, \dots, X_n$ sont iid de fonction de répartition continue $F$ entièrement
spécifiée utilise $D_n = \sup_x|F_n(x) - F(x)|$, $F_n$ étant la fonction de répartition empirique. Sous l'\omterm{def:qm:testing-and-multiple-testing:test}{hypothèse nulle}, $\P(\sqrt n\,D_n > \lambda)
\to 2\sum_{k \ge 1}(-1)^{k-1}e^{-2k^2\lambda^2}$, quelle que soit $F$~; la valeur critique à 5~\% est $\lambda = 1.358$.
\end{definition}

Mille rendements quotidiens tirés d'une loi $t$ de Student à quatre degrés de liberté, ramenés à une volatilité de 1~\%, donnent $D_n = 0.0615$
contre la loi normale de même volatilité~: au-dessus de la valeur critique $1.358/\sqrt{1\,000} = 0.0429$, soit une $p$-valeur de
0.0010. Deux mises en garde accompagnent le test. Il est le moins sensible dans les queues, où $F_n - F$ reste petit même lorsque les queues sont mal spécifiées~;
et lorsque la moyenne et la variance sont estimées à partir des mêmes données, $D_n$ est plus petit que ne le suppose la loi tabulée, de sorte que les
valeurs critiques doivent être simulées pour la famille ajustée (la correction de Lilliefors).

La question de la durée d'attente est l'usage le plus pointu du niveau et de la puissance dans un groupe de recherche.

\begin{proposition}[Taille d'échantillon pour un ratio de Sharpe]\label{prop:qm:testing-and-multiple-testing:years}
Avec $T$ années de rendements indépendants et un \omterm{def:qm:estimation:sharpe}{ratio de Sharpe} annuel $\mathrm{SR}$, le test unilatéral de $\mathrm{SR} = 0$ au niveau
$\alpha$ a une puissance approximativement égale à $\Phi\bigl(\mathrm{SR}\sqrt T - z_{1-\alpha}\bigr)$, et atteint la puissance $1 - \beta$ après
\[
T = \Bigl(\frac{z_{1-\alpha} + z_{1-\beta}}{\mathrm{SR}}\Bigr)^2 \text{ années}.
\]
\end{proposition}

\begin{proof}
Par la \omterm{met:qm:estimation:delta}{méthode delta} du chapitre~11, l'estimation annualisée suit approximativement $\mathcal N(\mathrm{SR}, 1/T)$ lorsque le ratio quotidien est
petit. Le test rejette lorsque $\widehat{\mathrm{SR}}\sqrt T > z_{1-\alpha}$, ce qui a pour probabilité $\Phi(\mathrm{SR}\sqrt T - z_{1-\alpha})$~;
en posant celle-ci égale à $1 - \beta$, on obtient $T$.
\end{proof}

Avec un niveau de 5~\% et une puissance de 80~\%, un \omterm{def:qm:estimation:sharpe}{ratio de Sharpe} de 2 exige 1.5 année, un \omterm{def:qm:estimation:sharpe}{ratio de Sharpe} de 1 exige 6.2 années et un \omterm{def:qm:estimation:sharpe}{ratio de Sharpe} de 0.5
exige 24.7 années (la \cref{fig:qm:testing-and-multiple-testing:power}). Cinq années d'un \omterm{def:qm:estimation:sharpe}{ratio de Sharpe} véritable de 1 sont détectées dans 72~\% des cas~: plus
d'une bonne stratégie sur quatre échoue à son premier test honnête.

\begin{omfigure}
\begin{tikzpicture}
\begin{axis}[width=10cm,height=5cm,xlabel={années de rendements quotidiens},ylabel={puissance à un niveau de 5~\%},
  tick label style={font=\small},label style={font=\small},xmin=0,xmax=15,ymin=0,ymax=1,grid=major,grid style={gray!25},
  legend columns=4,legend style={font=\footnotesize,at={(0.5,-0.32)},anchor=north,draw=none,fill=none,
  /tikz/every even column/.append style={column sep=8pt}},table/col sep=comma]
\addplot[omDef,very thick] table[x=years,y=sr2]{figdata/methods/12-testing-and-multiple-testing/power.csv};
\addplot[omThm,very thick,dashed] table[x=years,y=sr1]{figdata/methods/12-testing-and-multiple-testing/power.csv};
\addplot[black!70,very thick,dotted] table[x=years,y=sr05]{figdata/methods/12-testing-and-multiple-testing/power.csv};
\addplot[black!50,thin,sharp plot] coordinates {(0,0.8) (15,0.8)};
\legend{$\mathrm{SR} = 2$,$\mathrm{SR} = 1$,$\mathrm{SR} = 0.5$,80~\% de puissance}
\end{axis}
\end{tikzpicture}

\omcaption{Puissance du test unilatéral à 5~\% d'un \omterm{def:qm:estimation:sharpe}{ratio de Sharpe} nul, en fonction des années de rendements quotidiens indépendants, pour des \omterm{def:qm:estimation:sharpe}{ratios de Sharpe} annuels
véritables de 2, 1 et 0.5. La courbe à 80~\% de puissance est franchie après 1.5, 6.2 et 24.7 années. Données~: le module de
tutoriel du chapitre.}\label{fig:qm:testing-and-multiple-testing:power}
\end{omfigure}

\section{Le jardin aux sentiers qui bifurquent}

\begin{definition}[Tests multiples, data snooping, jardin aux sentiers qui bifurquent]\label{def:qm:testing-and-multiple-testing:snooping}
Les \emph{tests multiples}\index{tests multiples} consistent à effectuer de nombreux tests et à en rapporter certains. Le \emph{data
snooping}\index{data snooping} consiste à choisir un modèle, une stratégie ou un paramètre en cherchant sur les mêmes données qui serviront ensuite à
les tester, de sorte que la statistique rapportée est la meilleure d'une recherche présentée comme un test unique. Le \emph{jardin aux sentiers qui
bifurquent}\index{jardin aux sentiers qui bifurquent} est la version implicite~: chaque choix que l'analyste aurait fait autrement si les données
avaient été autres (la période d'échantillon, l'univers, le filtre, le modèle de coûts) multiplie les tests réellement effectués, même lorsque
une seule analyse est menée.
\end{definition}

Le nom est celui de Gelman et Loken (2013), et le propos est le leur~: un chercheur qui ne lance jamais deux backtests peut encore faire du \omterm{def:qm:testing-and-multiple-testing:snooping}{data snooping},
parce que le seul backtest a été façonné par l'examen des données. Une recherche n'est corrigée que si elle est comptée.

\begin{proposition}[Maximum de statistiques indépendantes]\label{prop:qm:testing-and-multiple-testing:max}
Si $Z_1, \dots, Z_N$ sont des normales centrées réduites indépendantes, $\P(\max_i Z_i \ge c) = 1 - \Phi(c)^N$, et $\E[\max_i Z_i] \sim
\sqrt{2\ln N}$ lorsque $N \to \infty$.
\end{proposition}

\begin{proof}
$\P(\max_i Z_i < c) = \prod_i\P(Z_i < c) = \Phi(c)^N$. Pour la croissance, $1 - \Phi(c) \approx \varphi(c)/c$, et le maximum se situe là où
$N(1 - \Phi(c)) \approx 1$, soit $c^2/2 \approx \ln N$ au premier ordre.
\end{proof}

La convergence est lente~: pour $N = 200$, le maximum espéré vaut 2.75, et non $\sqrt{2\ln200} = 3.26$. Ce maximum de deux cents
statistiques de bruit franchit presque toujours le seuil conventionnel $t > 2$, et le 3.1 de l'accroche environ une fois sur six
(\cref{fig:qm:testing-and-multiple-testing:maxt}). Un test isolé à $t = 3.1$ a $p = 0.001$~; le meilleur de deux cents tests indépendants à $t = 3.1$
a $p = 0.176$.

\begin{omfigure}
\begin{tikzpicture}
\begin{semilogyaxis}[width=10cm,height=5.2cm,xlabel={seuil $c$},ylabel={$\P(\max t \ge c)$ sous l'hypothèse nulle},
  tick label style={font=\small},label style={font=\small},xmin=0,xmax=4.5,ymin=1e-5,ymax=1.5,grid=major,grid style={gray!25},
  legend columns=3,legend style={font=\footnotesize,at={(0.5,-0.32)},anchor=north,draw=none,fill=none,
  /tikz/every even column/.append style={column sep=8pt}},table/col sep=comma]
\addplot[black!70,very thick,dotted] table[x=c,y=single]{figdata/methods/12-testing-and-multiple-testing/maxt.csv};
\addplot[omDef,very thick] table[x=c,y=correlated]{figdata/methods/12-testing-and-multiple-testing/maxt.csv};
\addplot[omThm,very thick,dashed] table[x=c,y=independent]{figdata/methods/12-testing-and-multiple-testing/maxt.csv};
\addplot[black!50,thin,sharp plot] coordinates {(3.11,1e-5) (3.11,1.5)};
\legend{un test,{200 variantes corrélées},{200 tests indépendants}}
\end{semilogyaxis}
\end{tikzpicture}

\omcaption{Probabilité que la plus grande statistique $t$ atteigne $c$ lorsque rien ne fonctionne~: un test, les deux cents variantes de momentum
du tutoriel (simulation gaussienne avec leur corrélation estimée), et deux cents tests indépendants. À la
valeur 3.11 de la meilleure variante (trait vertical)~: 0.0009, 0.030 et 0.17. Données~: le tutoriel du chapitre, avec graine fixée.}\label{fig:qm:testing-and-multiple-testing:maxt}
\end{omfigure}

Les recherches réelles sont corrélées. La famille du tutoriel compte deux cents règles de tendance qui retiennent le signe du
rendement glissant de $L$ jours, pour des horizons $L = 5, \dots, 204$ jours, appliquées à dix ans d'une marche aléatoire sans dérive, à 1~\% de volatilité quotidienne~: aucune règle
n'a le moindre edge. Des horizons voisins partagent presque tous leurs signaux (le P\&L quotidien de la meilleure variante a une corrélation de 0.95 avec celui de son
voisin), les horizons éloignés beaucoup moins (0.28 avec la règle à 5 jours, 0.20 avec celle à 204 jours). Dans l'historique présenté, 33 des
200 variantes ont une $p$-valeur unilatérale inférieure à 5~\%, la moyenne des statistiques $t$ sur la famille vaut 1.04 au lieu de zéro, car
chaque variante voit les mêmes épisodes de tendance, et la meilleure, avec un horizon de 42 jours, a $t = 3.11$ et un \omterm{def:qm:estimation:sharpe}{ratio de Sharpe} annualisé de
0.98 (\cref{fig:qm:testing-and-multiple-testing:variants}). Cet historique est le 336\up{e} que nous avons simulé~: nous avons tiré des historiques jusqu'à ce que la meilleure variante s'approche
du 3.1 de l'accroche, ce qui met en pratique, sur le chapitre lui-même, son propre sujet.

\begin{omfigure}
\begin{tikzpicture}
\begin{axis}[width=11cm,height=5cm,xlabel={horizon $L$ (jours)},ylabel={statistique $t$},
  tick label style={font=\small},label style={font=\small},xmin=0,xmax=210,ymin=-1.5,ymax=4,grid=major,grid style={gray!25},
  legend columns=4,legend style={font=\footnotesize,at={(0.5,-0.32)},anchor=north,draw=none,fill=none,
  /tikz/every even column/.append style={column sep=6pt}},table/col sep=comma]
\addplot[omDef,thick,mark=*,mark size=0.8pt] table[x=lookback,y=t]{figdata/methods/12-testing-and-multiple-testing/variants.csv};
\addplot[black!60,thin,dashed,sharp plot] coordinates {(0,1.645) (210,1.645)};
\addplot[omThm,very thick,dashdotted,sharp plot] coordinates {(0,2.916) (210,2.916)};
\addplot[black,thick,dotted,sharp plot] coordinates {(0,3.481) (210,3.481)};
\legend{variantes,un test à 5~\%,max-$t$ à 5~\%,Bonferroni à 5~\%}
\end{axis}
\end{tikzpicture}

\omcaption{Les statistiques $t$ des deux cents règles de tendance sur une marche aléatoire simulée de dix ans, en fonction de leur horizon, avec
trois seuils unilatéraux à 5~\%~: un test isolé (1.645), la valeur critique de la statistique du maximum pour cette famille corrélée (2.92) et
celle de Bonferroni (3.48). La meilleure règle, à 42 jours, a $t = 3.11$. Données~: le tutoriel du chapitre, avec graine fixée.}\label{fig:qm:testing-and-multiple-testing:variants}
\end{omfigure}

\section{Erreur familiale et fausses découvertes}

\begin{definition}[Taux d'erreur familial, correction de Bonferroni, procédure de Holm]\label{def:qm:testing-and-multiple-testing:fwer}
Pour $m$ hypothèses dont $m_0$ sont vraies, le \emph{taux d'erreur familial}\index{taux d'erreur familial} (FWER) d'une procédure est
la probabilité qu'elle rejette au moins une \omterm{def:qm:testing-and-multiple-testing:test}{hypothèse nulle} vraie. La \emph{correction de Bonferroni}\index{correction de Bonferroni} rejette
$H_i$ lorsque $p_i \le \alpha/m$. La \emph{procédure de Holm}\index{procédure de Holm} ordonne les $p$-valeurs $p_{(1)} \le \dots \le p_{(m)}$ et
rejette $H_{(1)}, \dots, H_{(k)}$, où $k$ est le dernier indice tel que $p_{(j)} \le \alpha/(m - j + 1)$ pour tout $j \le k$.
\end{definition}

\begin{proposition}[Bonferroni et Holm contrôlent le taux d'erreur familial]\label{prop:qm:testing-and-multiple-testing:fwer}
Quelle que soit la dépendance entre les tests, les deux procédures ont un FWER au plus égal à $\alpha$~; Holm rejette tout ce que rejette Bonferroni.
\end{proposition}

\begin{proof}
Pour Bonferroni~: l'inégalité de Boole appliquée aux \omterm{def:qm:testing-and-multiple-testing:test}{hypothèses nulles} vraies donne $\P(\exists i \in \mathcal H_0: p_i \le \alpha/m) \le m_0\alpha/m \le \alpha$.
Pour Holm~: soit $j^*$ le rang de la plus petite $p$-valeur parmi les \omterm{def:qm:testing-and-multiple-testing:test}{hypothèses nulles} vraies. Au plus $m - m_0$ hypothèses la précèdent, donc $j^* \le
m - m_0 + 1$, et Holm ne rejette une \omterm{def:qm:testing-and-multiple-testing:test}{hypothèse nulle} vraie que si $p_{(j^*)} \le \alpha/(m - j^* + 1) \le \alpha/m_0$, ce qui, par l'inégalité de Boole sur
les $m_0$ \omterm{def:qm:testing-and-multiple-testing:test}{hypothèses nulles} vraies, a une probabilité au plus $\alpha$. Les seuils de Holm ne sont jamais inférieurs à $\alpha/m$.
\end{proof}

Maîtriser la probabilité de commettre la moindre fausse découverte est le bon objectif lorsqu'une fausse découverte coûte cher, par exemple quand une stratégie
unique passe en production. Quand un groupe de recherche passe au crible des centaines de signaux pour un modèle combiné, quelques fausses
découvertes parmi de nombreuses vraies coûtent peu, et n'en exiger aucune revient à jeter la plupart des vraies.

\begin{definition}[Taux de fausses découvertes, procédure de Benjamini--Hochberg]\label{def:qm:testing-and-multiple-testing:fdr}
Si une procédure effectue $R$ rejets dont $V$ sont des \omterm{def:qm:testing-and-multiple-testing:test}{hypothèses nulles} vraies, son \emph{taux de fausses découvertes}\index{taux de fausses découvertes} est
$\E[V/\max(R, 1)]$. La \emph{procédure de Benjamini--Hochberg}\index{procédure de Benjamini--Hochberg} au niveau $q$ rejette $H_{(1)},
\dots, H_{(k)}$ pour le plus grand $k$ tel que $p_{(k)} \le kq/m$.
\end{definition}

\begin{theorem}[Benjamini--Hochberg]\label{thm:qm:testing-and-multiple-testing:bh}
Si les $p$-valeurs des \omterm{def:qm:testing-and-multiple-testing:test}{hypothèses nulles} vraies sont indépendantes entre elles et des autres, la \omterm{def:qm:testing-and-multiple-testing:fdr}{procédure de Benjamini--Hochberg} a un \omterm{def:qm:testing-and-multiple-testing:fdr}{taux de fausses
découvertes} exactement égal à $q\,m_0/m \le q$. Il reste inférieur ou égal à $q$ sous une dépendance de régression positive, et au plus $q$ sous toute
dépendance si $q$ est remplacé par $q/\sum_{k=1}^m k^{-1}$ (Benjamini et Yekutieli, 2001).
\end{theorem}

\begin{example}[Cinq $p$-valeurs]\label{ex:qm:testing-and-multiple-testing:five}
Prenons $p = (0.005, 0.01, 0.03, 0.04, 0.2)$ et $\alpha = q = 5\%$. Le seuil de Bonferroni, $0.01$, en rejette deux. Holm les compare dans
l'ordre avec $0.05/5, 0.05/4, 0.05/3, \dots$~: $0.005 \le 0.01$ et $0.01 \le 0.0125$, puis $0.03 > 0.0167$ l'arrête à deux. Benjamini--Hochberg
compare avec $0.01, 0.02, 0.03, 0.04, 0.05$~: le plus grand $k$ tel que $p_{(k)} \le 0.01k$ est $k = 4$, donc elle en rejette quatre. Sous forme de
$p$-valeurs ajustées (le plus petit niveau auquel chacune est rejetée)~: Bonferroni $(0.025, 0.05, 0.15, 0.2, 1)$, Holm $(0.025, 0.04, 0.09, 0.09, 0.2)$,
Benjamini--Hochberg $(0.025, 0.025, 0.05, 0.05, 0.2)$.
\end{example}

Le compromis se voit à grande échelle. Parmi mille signaux candidats dont cent sont réels (avec des statistiques $t$ centrées sur 3), le
test naïf à 5~\% produit 136 découvertes dont 45 sont fausses~; Bonferroni et Holm en font environ 18, avec une fausse découverte dans 5~\% des
criblages~; Benjamini--Hochberg à 5~\% en produit 63, dont 2.9 sont fausses en moyenne, soit un \omterm{def:qm:testing-and-multiple-testing:fdr}{taux de fausses découvertes} de 4.4~\%, à comparer avec
celui du théorème, $0.05 \times 900/1\,000 = 4.5\%$ (\cref{fig:qm:testing-and-multiple-testing:fdr}). Il paie sa puissance par un \omterm{def:qm:testing-and-multiple-testing:fwer}{taux d'erreur familial} de 92~\%~:
presque chaque criblage contient un faux signal, et c'est là le choix de conception.

\begin{omfigure}
\begin{tikzpicture}
\begin{axis}[width=10cm,height=4.8cm,xbar stacked,bar width=10pt,
  xlabel={découvertes par criblage (moyenne sur 500)},
  ytick={0,1,2,3},yticklabels={aucune correction,Bonferroni,Holm,Benjamini--Hochberg},
  tick label style={font=\small},label style={font=\small},
  xmin=0,xmax=150,ymin=-0.6,ymax=3.6,grid=major,grid style={gray!25},
  legend columns=2,legend style={font=\footnotesize,at={(0.5,-0.4)},anchor=north,draw=none,fill=none,
  /tikz/every even column/.append style={column sep=8pt}},table/col sep=comma]
\addplot[fill=omDef!60,draw=omDef] table[x=true,y=k]{figdata/methods/12-testing-and-multiple-testing/fdr.csv};
\addplot[fill=omThm!50,draw=omThm] table[x=false,y=k]{figdata/methods/12-testing-and-multiple-testing/fdr.csv};
\legend{découvertes vraies,fausses découvertes}
\end{axis}
\end{tikzpicture}

\omcaption{Mille signaux, dont cent sont réels~: moyenne des découvertes vraies et fausses de quatre règles à 5~\%. Aucune correction~:
91 vraies, 45 fausses. Bonferroni et Holm~: 18 vraies, 0.05 fausse. Benjamini--Hochberg~: 61 vraies, 2.9 fausses (un \omterm{def:qm:testing-and-multiple-testing:fdr}{taux de fausses
découvertes} de 4.4~\%). Données~: le tutoriel du chapitre, avec graine fixée.}\label{fig:qm:testing-and-multiple-testing:fdr}
\end{omfigure}

Sur la famille des signaux de momentum, aucune des deux corrections n'est clémente. Bonferroni et Holm donnent à la meilleure variante une $p$-valeur ajustée de $200
\times 0.00094 = 0.187$~; Benjamini--Hochberg donne 0.076, plus petite parce que sa procédure ascendante emprunte de la force aux 33 autres petites
$p$-valeurs. Aucune variante ne survit à 5~\%. Mais ces deux corrections ont été construites pour le pire cas de dépendance, et deux cents
variantes qui partagent l'essentiel de leurs signaux sont loin de deux cents chances indépendantes.

\section{Tests de contrôle de réalité pour la recherche de stratégies}

\begin{definition}[Reality Check, test de capacité prédictive supérieure]\label{def:qm:testing-and-multiple-testing:rc}
Pour $K$ stratégies de séries de performance $d_{k,t}$ en excès d'un indice de référence, le \emph{Reality Check}\index{Reality Check}
(White, 2000) teste $H_0: \max_k\E[d_k] \le 0$ à l'aide de la statistique $\max_k\sqrt n\,\bar d_k$, dont la loi nulle est estimée par
le bootstrap de $\max_k\sqrt n(\bar d_k^* - \bar d_k)$. Le \emph{test de capacité prédictive supérieure}\index{test de capacité prédictive supérieure}
(Hansen, 2005) studentise chaque $\bar d_k$ et ne recentre que les stratégies qui ne sont pas nettement moins bonnes que l'indice de référence, de sorte
que l'ajout de stratégies médiocres à la recherche ne rend pas le test plus conservateur.
\end{definition}

Les deux tests remplacent la borne de l'union par la loi effective du maximum, et le bootstrap (chapitre~13) préserve la dépendance entre
stratégies et dans le temps. Lorsque les statistiques sont approximativement conjointement gaussiennes, une simulation gaussienne avec
leur corrélation estimée produit le même résultat, sans coût supplémentaire.

\begin{method}[$p$-valeurs de la statistique maximale]\label{met:qm:testing-and-multiple-testing:maxt}
Étant donné des statistiques $t$ $t_1, \dots, t_m$ et la matrice de corrélation $R$ des séries de P\&L sous-jacentes~: tirer $Z^{(s)} \sim \mathcal N(0, R)$
pour $s = 1, \dots, S$~; la $p$-valeur ajustée de $H_i$ est la part des tirages pour lesquels $\max_j Z_j^{(s)} \ge t_i$. La version
descendante (Romano et Wolf, 2005) prend le maximum, pour la $k$-ième statistique la plus grande, seulement sur les hypothèses non encore
rejetées. Les deux contrôlent le \omterm{def:qm:testing-and-multiple-testing:fwer}{taux d'erreur familial}~; le nombre effectif d'essais indépendants de la recherche est le $N$ qui résout $1 - (1 -
p_{\text{naïve}})^N = p_{\text{ajustée}}$.
\end{method}

Pour la famille de momentum, 200\,000 tirages à partir de la corrélation estimée donnent à la meilleure variante une $p$-valeur ajustée de 0.030,
contre 0.187 pour Bonferroni et 0.171 pour deux cents tests indépendants. La recherche sur deux cents horizons corrélés revient
à 32 essais indépendants. La valeur critique à 5~\% de la statistique maximale est 2.92, nettement en dessous de 3.48 pour Bonferroni, et trois variantes
la franchissent. Le contrôle qui compte~: sur 2\,000 historiques décennaux neufs, la meilleure variante atteint $t = 3.11$ dans 2.6~\% des
cas, proche des 3.0~\% simulés.

Ainsi, le test corrigé rejette, à 5~\%, une stratégie qui n'a aucun avantage. Il en a le droit, 3~\% du temps~; et cet historique était le
336\up{e} que nous avons tiré. La correction tient compte des deux cents horizons dont elle a été informée, et de rien d'autre~: ni des
historiques écartés, ni des autres familles de signaux que l'équipe a testées l'an dernier. Harvey, Liu et Zhu (2016) aboutissent à une
conclusion voisine pour les facteurs actions publiés, en soutenant qu'un nouveau facteur devrait franchir une statistique $t$ d'environ 3 plutôt que 2.

\begin{definition}[Ratio de Sharpe déflaté]\label{def:qm:testing-and-multiple-testing:dsr}
Pour un \omterm{def:qm:estimation:sharpe}{ratio de Sharpe} $\widehat{\mathrm{SR}}$ (par période) estimé à partir de $n$ rendements, d'asymétrie $\gamma_3$ et d'aplatissement
$\gamma_4$, la probabilité que le ratio vrai dépasse $\mathrm{SR}_0$ vaut approximativement
\[
\Phi\Bigl((\widehat{\mathrm{SR}} - \mathrm{SR}_0)\sqrt{n - 1}\Big/\sqrt{1 - \gamma_3\widehat{\mathrm{SR}} + \tfrac{\gamma_4 - 1}4\widehat{\mathrm{SR}}^2}\Bigr).
\]
Le \emph{ratio de Sharpe déflaté}\index{ratio de Sharpe déflaté} (Bailey et López de Prado, 2014) est cette probabilité, où $\mathrm{SR}_0$
est l'espérance du maximum de $N$ \omterm{def:qm:estimation:sharpe}{ratios de Sharpe} nuls de variance $V$~: $\mathrm{SR}_0 = \sqrt V\bigl((1 - \gamma_E)\Phi^{-1}(1 - 1/N) +
\gamma_E\Phi^{-1}(1 - 1/(Ne))\bigr)$, $\gamma_E$ étant la constante d'Euler--Mascheroni.
\end{definition}

Pour la meilleure variante de momentum, la probabilité que son \omterm{def:qm:estimation:sharpe}{ratio de Sharpe} vrai soit positif vaut 0.999, selon la vue à test unique.
Déflatée pour deux cents essais (avec $V = 1/n$, la variance nulle), la référence devient un 0.87 annualisé et le \omterm{def:qm:testing-and-multiple-testing:dsr}{ratio de Sharpe déflaté} 0.63~;
pour les 32 essais effectifs, 0.67 et 0.84. Aucun n'atteint les 0.95 qu'exigerait un desk, et le nombre d'essais, que le chercheur
contrôle et que le lecteur voit rarement, pèse davantage sur la réponse que les données.

\section{Tutoriel~: la meilleure parmi deux cents}\label{tut:qm:testing-and-multiple-testing:best}

\begin{tutorial}
\textbf{Objectif.} Parcourir deux cents règles de tendance sans valeur, retenir la meilleure, et mesurer son caractère surprenant selon cinq corrections.
\textbf{État final~:} les \cref{fig:qm:testing-and-multiple-testing:variants,fig:qm:testing-and-multiple-testing:maxt} et les $p$-valeurs ajustées 0.187, 0.076 et 0.030.
\begin{tutsteps}
\item \textbf{$p$-valeurs ajustées} par la procédure descendante de Holm et la procédure ascendante de Benjamini--Hochberg.
\omcode{code/firm/multitest/firm_multitest.py}{68}{88}{$p$-valeurs ajustées de Holm et de Benjamini--Hochberg.}\label{lst:qm:testing-and-multiple-testing:holm}
\item \textbf{$p$-valeurs de la statistique maximale} à partir d'une \omterm{def:qm:testing-and-multiple-testing:test}{hypothèse nulle} gaussienne avec la corrélation estimée, ou de toute matrice de tirages nuls.
\omcode{code/firm/multitest/firm_multitest.py}{97}{133}{$p$-valeurs de la statistique maximale, en une étape et en descente.}\label{lst:qm:testing-and-multiple-testing:maxt}
\item \textbf{Exécuter} \texttt{problem()} dans \texttt{qm\_testing.py}, puis \texttt{family\_exceedance(3.11)} pour vérifier la $p$-valeur simulée sur
des historiques neufs, et \texttt{fig\_testing.py}.
\end{tutsteps}
\textbf{Pour aller plus loin.} Ajouter les règles de renversement symétriques et observer croître le nombre effectif d'essais~; remplacer l'\omterm{def:qm:testing-and-multiple-testing:test}{hypothèse nulle}
gaussienne par le bootstrap stationnaire du chapitre~13 via le point d'accroche \texttt{null\_draws}.
\end{tutorial}

\section{Construction~: le module de tests multiples}\label{bld:qm:testing-and-multiple-testing:multitest}

\begin{build}
\bfield{Rôle} Chaque stratégie que la petite société promeut depuis la recherche porte avec elle la taille de la recherche dont elle est issue et une
$p$-valeur ajustée~; ce module les calcule.
\bfield{Interface} \texttt{bonferroni(p)}, \texttt{holm(p)}, \texttt{benjamini\_hochberg(p)}, \texttt{benjamini\_yekutieli(p)} renvoient des
$p$-valeurs ajustées~; \texttt{maxt\_pvalues(t, corr, n\_sim, seed, null\_draws, stepdown)}~; \texttt{effective\_trials(p\_single,
p\_family)}~; \texttt{expected\_max\_sr(n\_trials, sr\_var)}~; \texttt{deflated\_sharpe(sr, n\_obs, skew, kurt, sr0)}~; \texttt{norm\_cdf},
\texttt{norm\_ppf}.
\bfield{Règles} $p$-valeurs ajustées et non indicateurs de rejet, afin que le niveau soit choisi par l'appelant~; statistiques unilatérales pour la
recherche de stratégies~; les tirages nuls sont un paramètre, de sorte que le bootstrap remplace la gaussienne sans toucher aux appelants.
\bfield{Tests d'acceptation} \texttt{code/firm/multitest/tests/}~: l'exemple des cinq $p$-valeurs~; contrôle du \omterm{def:qm:testing-and-multiple-testing:fwer}{taux d'erreur familial} et du \omterm{def:qm:testing-and-multiple-testing:fdr}{taux de fausses découvertes} par simulation~; la
$p$-valeur de la statistique maximale égale la formule indépendante pour $R = I$ et la $p$-valeur naïve pour des tests parfaitement corrélés~;
l'approximation de l'espérance du maximum à 1.5~\% près.
\bfield{Pour aller plus loin} Le test SPA de Hansen avec le bootstrap stationnaire~; la procédure descendante de Romano--Wolf avec des tirages nuls bootstrap~; la
probabilité de surapprentissage du backtest par validation croisée combinatoire.
\end{build}

\begin{omsources}
\item J.~Neyman et E.~S.~Pearson, ``On the problem of the most efficient tests of statistical hypotheses'', \emph{Philosophical
Transactions of the Royal Society A} 231, 1933.
\item S.~Holm, ``A simple sequentially rejective multiple test procedure'', \emph{Scandinavian Journal of Statistics} 6, 1979.
\item Y.~Benjamini et Y.~Hochberg, ``Controlling the false discovery rate: a practical and powerful approach to multiple testing'',
\emph{Journal of the Royal Statistical Society B} 57, 1995.
\item Y.~Benjamini et D.~Yekutieli, ``The control of the false discovery rate in multiple testing under dependency'', \emph{Annals of
Statistics} 29, 2001.
\item H.~White, ``A reality check for data snooping'', \emph{Econometrica} 68, 2000.
\item P.~R.~Hansen, ``A test for superior predictive ability'', \emph{Journal of Business and Economic Statistics} 23, 2005.
\item J.~P.~Romano et M.~Wolf, ``Stepwise multiple testing as formalized data snooping'', \emph{Econometrica} 73, 2005.
\item C.~R.~Harvey, Y.~Liu et H.~Zhu, ``\dots and the cross-section of expected returns'', \emph{Review of Financial Studies} 29, 2016.
\item D.~H.~Bailey et M.~López de Prado, ``The deflated Sharpe ratio'', \emph{Journal of Portfolio Management} 40(5), 2014.
\item A.~Gelman et E.~Loken, ``The garden of forking paths'', document de travail, 2013.
\end{omsources}

\section{Exercices}

\begin{exercise}[$\star$]\label{exo:qm:testing-and-multiple-testing:1}
Un backtest affiche $t = 2.5$. Calculer sa $p$-valeur unilatérale et sa $p$-valeur ajustée par la \omterm{def:qm:testing-and-multiple-testing:fwer}{correction de Bonferroni}
s'il est le meilleur de vingt variantes~?
\end{exercise}

\begin{exercise}[$\star$]\label{exo:qm:testing-and-multiple-testing:2}
Un test rejette lorsque $Z > 1.645$. Sous l'alternative d'intérêt $Z \sim \mathcal N(2, 1)$, calculer son niveau et sa puissance.
\end{exercise}

\begin{exercise}[$\star$]\label{exo:qm:testing-and-multiple-testing:3}
Appliquer la \omterm{def:qm:testing-and-multiple-testing:fwer}{correction de Bonferroni} et la \omterm{def:qm:testing-and-multiple-testing:fwer}{procédure de Holm} au niveau de 5~\% à $p = (0.001, 0.012, 0.02, 0.3)$.
\end{exercise}

\begin{exercise}[$\star\star$]\label{exo:qm:testing-and-multiple-testing:4}
Calculer le nombre d'années de rendements indépendants nécessaire à un test de niveau 5~\% pour détecter un \omterm{def:qm:estimation:sharpe}{ratio de Sharpe} annualisé de 0.7 avec une puissance de 90~\%.
\end{exercise}

\begin{exercise}[$\star\star$]\label{exo:qm:testing-and-multiple-testing:5}
Les rendements quotidiens sont gaussiens, de moyenne nulle. Les 500 premiers jours ont une variance empirique de $(1.0\%)^2$, les 500 suivants $(1.2\%)^2$. Calculer la
statistique du \omterm{def:qm:testing-and-multiple-testing:lr}{test du rapport de vraisemblance} pour une variance commune et sa $p$-valeur.
\end{exercise}

\begin{exercise}[$\star\star$]\label{exo:qm:testing-and-multiple-testing:6}
Pour dix stratégies indépendantes de pur bruit, calculer la probabilité que la meilleure ait $t \ge 2$ et le maximum
espéré selon la formule du \omterm{def:qm:testing-and-multiple-testing:dsr}{ratio de Sharpe déflaté}.
\end{exercise}

\begin{exercise}[$\star\star\star$]\label{exo:qm:testing-and-multiple-testing:7}
\emph{Programmation.} Reprendre le criblage de mille signaux avec des statistiques nulles partageant un facteur commun (corrélation 0.5). Vérifier si
la \omterm{def:qm:testing-and-multiple-testing:fdr}{procédure de Benjamini--Hochberg} contrôle encore le \omterm{def:qm:testing-and-multiple-testing:fdr}{taux de fausses découvertes} à 5~\%, et calculer la fréquence à laquelle
la proportion réalisée de fausses découvertes dépasse 10~\%.
\end{exercise}

\begin{exercise}[$\star\star\star$]\label{exo:qm:testing-and-multiple-testing:8}
\emph{Trouver la faille.} «~Nous avons testé notre signal séparément sur douze marchés. Il est significatif à 5~\% sur trois d'entre eux, où nous le
négocions désormais~; les trois backtests constituent notre estimation de sa performance.~»
\end{exercise}

\section{Problème~: deux cents variantes de momentum}

\begin{problem}[{Problème du week-end --- à quel point le meilleur résultat d'une recherche corrélée est-il surprenant~?}]\label{pb:qm:testing-and-multiple-testing:1}
Une équipe réalise le backtest de deux cents règles de tendance (conserver le signe du rendement glissant sur $L$ jours, $L = 5, \dots, 204$) sur dix ans d'un
marché qui, à leur insu, est une marche aléatoire sans dérive, de volatilité quotidienne de 1~\%. L'historique est celui du chapitre, généré à partir d'une graine.

\medskip\noindent\textbf{Partie I --- La recherche.}
\begin{enumerate}
\item Identifier la fenêtre d'observation gagnante, sa statistique $t$ et son \omterm{def:qm:estimation:sharpe}{ratio de Sharpe} annualisé.
\item Calculer sa $p$-valeur unilatérale naïve.
\item Dénombrer les règles, parmi les deux cents, qui passent un test unilatéral au niveau de 5~\%.
\item Calculer la moyenne des statistiques $t$ sur la famille et expliquer pourquoi elle n'est pas proche de zéro.
\item Calculer la corrélation entre le P\&L de la meilleure règle et celui de sa voisine, puis avec celui de la règle à 5 jours et celui de la règle à 204 jours.
\end{enumerate}

\medskip\noindent\textbf{Partie II --- Corrections.}
\begin{enumerate}[resume]
\item Calculer la $p$-valeur ajustée par Bonferroni et la valeur critique de Bonferroni.
\item Donner le résultat de la \omterm{def:qm:testing-and-multiple-testing:fwer}{procédure de Holm} pour la meilleure règle, et expliquer pourquoi il est identique à celui de Bonferroni.
\item Donner le résultat de la \omterm{def:qm:testing-and-multiple-testing:fdr}{procédure de Benjamini--Hochberg}, et expliquer pourquoi il est plus petit.
\item Calculer la $p$-valeur ajustée qu'on obtiendrait si les deux cents règles étaient indépendantes.
\item Calculer, à partir de la corrélation estimée, la $p$-valeur et la valeur critique ajustées par la statistique du maximum.
\end{enumerate}

\medskip\noindent\textbf{Partie III --- Ce que vaut la recherche.}
\begin{enumerate}[resume]
\item Calculer le nombre effectif d'essais indépendants.
\item Dénombrer les variantes qui franchissent le seuil de la statistique du maximum.
\item Sur 2\,000 historiques nouveaux, calculer la fréquence à laquelle la meilleure règle atteint $t = 3.11$.
\item Calculer le \omterm{def:qm:estimation:sharpe}{ratio de Sharpe} probabiliste et le \omterm{def:qm:testing-and-multiple-testing:dsr}{ratio de Sharpe déflaté}, avec 200 essais puis avec le nombre effectif d'essais.
\item Vérifier si les réponses changent lorsque les rendements quotidiens ont des queues de Student-$t$ à quatre degrés de liberté.
\end{enumerate}

\medskip\noindent\textbf{Partie IV --- Jugement.}
\begin{enumerate}[resume]
\item Les règles n'ont aucun edge par construction. Expliquer pourquoi le test corrigé rejette l'\omterm{def:qm:testing-and-multiple-testing:test}{hypothèse nulle}.
\item Énumérer les conditions qu'il faudrait exiger avant de négocier la règle à 42 jours.
\item Indiquer ce que le rapport de recherche doit consigner sur la recherche.
\item Énoncer le \emph{résultat nommé}~: la $p$-valeur ajustée de la meilleure parmi deux cents variantes corrélées et le nombre effectif
d'essais.
\item En une phrase~: contre quoi une correction des \omterm{def:qm:testing-and-multiple-testing:snooping}{tests multiples} protège-t-elle, et contre quoi ne protège-t-elle pas~?
\end{enumerate}
\end{problem}

\section{Questions d'entretien}

\begin{interviewq}[$\star$ \iqroles{chercheur, trader}]\label{iq:qm:testing-and-multiple-testing:1}
Qu'est-ce qu'une $p$-valeur, et qu'est-ce qu'elle n'est pas~?
\end{interviewq}

\begin{interviewq}[$\star\star$ \iqroles{chercheur, trader}]\label{iq:qm:testing-and-multiple-testing:2}
Vous avez fait le backtest de cent stratégies, et la meilleure affiche $t = 2.8$. Est-elle réelle~?
\end{interviewq}

\begin{interviewq}[$\star\star$ \iqroles{chercheur, ingénieur ML}]\label{iq:qm:testing-and-multiple-testing:3}
Bonferroni ou Benjamini--Hochberg~: quand utiliseriez-vous l'un ou l'autre~?
\end{interviewq}

\begin{interviewq}[$\star\star$ \iqroles{chercheur, trader, risque}]\label{iq:qm:testing-and-multiple-testing:4}
Combien de temps une stratégie doit-elle fonctionner avant que vous puissiez confirmer un \omterm{def:qm:estimation:sharpe}{ratio de Sharpe} de 1~?
\end{interviewq}

\begin{interviewq}[$\star\star\star$ \iqroles{chercheur}]\label{iq:qm:testing-and-multiple-testing:5}
Vos deux cents variantes sont fortement corrélées. Comment corrigez-vous l'effet de la recherche~?
\end{interviewq}

\begin{interviewq}[$\star\star$ \iqroles{ingénieur ML, chercheur}]\label{iq:qm:testing-and-multiple-testing:6}
Qu'est-ce que le \omterm{def:qm:testing-and-multiple-testing:snooping}{jardin aux sentiers qui bifurquent}, et où se cache-t-il dans un pipeline d'apprentissage automatique~?
\end{interviewq}

\begin{interviewq}[$\star\star$ \iqroles{chercheur, risque}]\label{iq:qm:testing-and-multiple-testing:7}
Vous ajustez une loi aux rendements et un \omterm{def:qm:testing-and-multiple-testing:ks}{test de Kolmogorov--Smirnov} ne la rejette pas. Qu'avez-vous appris~?
\end{interviewq}
